% ---
% id: statistics/binomial-box/lem:widening
% kind: lem
% title: "Widening the box keeps its guarantees"
% summary: >-
%   Lowering the box's lower bounds and raising its upper ones gives a superset
%   that is failed no more often and has no larger infimum, and integer
%   endpoints so moved still meet their tail budgets. Moving them the other way
%   can break the tail budgets and the reported completeness figure ε*.
% prominence: standard
% status: accepted
% uses: [def:construction, def:endpoint-rule, def:interval, lem:guard, statistics/def:taxonomy, statistics/def:freq, statistics/def:envelope]
% mentions: [statistics/def:conservative-evaluation, prop:rule-exact]
% assumes: []
% refs: []
% provenance: newly-derived
% lean:
% ---

A \tref{conservative evaluation} (Definition~\ref{def:conservative-evaluation})
can only return a smaller lower and a larger upper endpoint than an exact one,
and the consuming engine's outward rounding of the emitted bounds moves them
the same way. The lemma says that every such move keeps what the construction
guarantees and costs only rate, and that the reverse move does not;
Proposition~\ref{prop:rule-exact}\,(iii) and Lemma~\ref{lem:guard}\,(ii) are
its counterparts for a rule's own interval.

\begin{lemma}\label{lem:widening}
Let $\acSet$ be the output of $\mathsf{Acc}_R$
(Definition~\ref{def:construction}) on an \tref{admissible input}
(Definition~\ref{def:endpoint-rule}) for an \tref{admissible endpoint rule}
$R$, with box bounds $\qlo,\qhi$ \eqref{eq:boxbounds} and
\tref{guarded endpoints} $k^{\mp}$ \eqref{eq:guard}. Let
$\tilde{\mathbf{q}}^{-},\tilde{\mathbf{q}}^{+}\in\mathbb{R}^{\Chatset}$
satisfy $\tilde q^{-}(\hat c)\le q^{-}(\hat c)$ and
$\tilde q^{+}(\hat c)\ge q^{+}(\hat c)$ for every $\hat c\in\Chatset$, and put
\begin{equation}\label{eq:widened-set}
  \acSet'\ \coloneqq\ \{\mathbf{q}\in\simplex{\Chatset}:\
  \tilde q^{-}(\hat c)\le q(\hat c)\le\tilde q^{+}(\hat c)\ \
  \forall\hat c\in\Chatset\}.
\end{equation}
\begin{enumerate}[label=(\roman*),leftmargin=2.2em]
\item \emph{(containment)} $\acSet\subseteq\acSet'$.
\item \emph{(failure probability)} For all $\Chatset$-valued random
variables $\Chat_1,\dots,\Chat_n$ on one probability space,
\begin{equation*}
  \Pr\big[\operatorname{freq}_{\Chat_1^n}\notin\acSet'\big]
  \ \le\ \Pr\big[\operatorname{freq}_{\Chat_1^n}\notin\acSet\big].
\end{equation*}
\item \emph{(infima)} For every function
$\Psi:\simplex{\Chatset}\to[-\infty,+\infty]$,
\begin{equation}\label{eq:widened-inf}
  \inf_{\mathbf{q}\in\acSet'}\Psi(\mathbf{q})
  \ \le\ \inf_{\mathbf{q}\in\acSet}\Psi(\mathbf{q}).
\end{equation}
\item \emph{(integer endpoints)} Let
$\tilde k^{-},\tilde k^{+}:\Chatcon\to\mathbb{Z}$ satisfy
$0\le\tilde k^{-}(\hat c)\le k^{-}(\hat c)$ and
$k^{+}(\hat c)\le\tilde k^{+}(\hat c)\le n$ for every $\hat c\in\Chatcon$,
and let $\tilde{\mathbf{q}}^{\mp}$ be given by \eqref{eq:boxbounds} with
$\tilde k^{\mp}$ in place of $k^{\mp}$. Then $\tilde{\mathbf{q}}^{\mp}$
satisfy the hypothesis above, so (i)--(iii) hold; the pair
$(\tilde{\mathbf{q}}^{-},\tilde{\mathbf{q}}^{+})$ meets the conditions of the
\tref{canonical box form} of Definition~\ref{def:taxonomy}, and its
denotation \eqref{eq:boxset} is $\acSet'$; and for every $\hat c\in\Chatcon$,
with $\Ncount{\hat c}^{\mp}$ as in Definition~\ref{def:endpoint-rule},
\begin{equation}\label{eq:widened-tails}
  \Pr\big[\Ncount{\hat c}^{-}<\tilde k^{-}(\hat c)\big]
  \le\Pr\big[\Ncount{\hat c}^{-}<k^{-}(\hat c)\big]
  \le\varepsilon^{-}(\hat c),
  \qquad
  \Pr\big[\Ncount{\hat c}^{+}>\tilde k^{+}(\hat c)\big]
  \le\Pr\big[\Ncount{\hat c}^{+}>k^{+}(\hat c)\big]
  \le\varepsilon^{+}(\hat c).
\end{equation}
Summed over $\Chatcon$, the tails at $\tilde k^{\mp}$ therefore total at most
the tails at $k^{\mp}$, and at most $\ecomp$.
\item \emph{(not the reverse)} Replacing $k^{-}(\hat c)$ by a larger integer
or $k^{+}(\hat c)$ by a smaller one keeps, in general, neither the budgets of
\eqref{eq:widened-tails} nor the tail sum at $k^{\mp}$ as a bound on the
probability that the \tref{frequency vector} leaves the set. Take
$\Chatset=\{\hat c_0,\hat c_1\}$, the \tref{degenerate envelope} at
$\honest=(\tfrac12,\tfrac12)$ (Definition~\ref{def:envelope}), $n=3$,
$\Chatcon=\{\hat c_0\}$,
$\varepsilon^{-}(\hat c_0)=\varepsilon^{+}(\hat c_0)=\tfrac25$,
$\ecomp=\tfrac45$ and the \tref{exact binomial quantile rule} of
Definition~\ref{def:interval}. Then $k^{-}(\hat c_0)=1$, $k^{+}(\hat c_0)=2$,
and the tails at these endpoints sum to $\tfrac14$. Replacing
$k^{-}(\hat c_0)$ by $2$, or $k^{+}(\hat c_0)$ by $1$, gives a tail of
$\tfrac12>\tfrac25$ on the replaced side; and for $\Chat_1,\Chat_2,\Chat_3$
independent with law $\honest$, the frequency vector lies outside the set
\eqref{eq:accset} built from the replaced endpoints with probability
$\tfrac58>\tfrac14$, while it lies outside $\acSet$ with probability
$\tfrac14$.
\end{enumerate}
\end{lemma}

\begin{proof}
\emph{(i).} Let $\mathbf{q}\in\acSet$ and $\hat c\in\Chatset$. By
\eqref{eq:accset} and the hypothesis,
$\tilde q^{-}(\hat c)\le q^{-}(\hat c)\le q(\hat c)\le q^{+}(\hat c)\le
\tilde q^{+}(\hat c)$. Since also $\mathbf{q}\in\simplex{\Chatset}$, every
defining condition of \eqref{eq:widened-set} holds, and
$\mathbf{q}\in\acSet'$.

\emph{(ii).} The frequency vector $\operatorname{freq}_{\Chat_1^n}$ of
Definition~\ref{def:freq} is a function of the finitely valued string
$\Chat_1^n$, so $\{\operatorname{freq}_{\Chat_1^n}\in S\}$ is the finite
union of the events $\{\Chat_1^n=\hat c_1^n\}$ over the strings with
$\operatorname{freq}_{\hat c_1^n}\in S$, an event for every set $S$. By (i),
$\{\operatorname{freq}_{\Chat_1^n}\in\acSet\}\subseteq
\{\operatorname{freq}_{\Chat_1^n}\in\acSet'\}$; taking complements,
$\{\operatorname{freq}_{\Chat_1^n}\notin\acSet'\}\subseteq
\{\operatorname{freq}_{\Chat_1^n}\notin\acSet\}$, and a probability measure
is monotone under inclusion.

\emph{(iii).} For every $\mathbf{q}\in\acSet$, (i) gives
$\mathbf{q}\in\acSet'$, so the infimum over $\acSet'$ is at most
$\Psi(\mathbf{q})$. The infimum over $\acSet'$ is therefore a lower bound of
$\{\Psi(\mathbf{q}):\mathbf{q}\in\acSet\}$, hence at most its greatest lower
bound, which is \eqref{eq:widened-inf}. This holds in the extended reals for
every $\Psi$, with no continuity or measurability required, and also when
$\acSet$ is empty, its infimum then being $+\infty$.

\emph{(iv).} \emph{The hypothesis.} For $\hat c\in\Chatcon$,
\eqref{eq:boxbounds} gives
$\tilde q^{-}(\hat c)=\tilde k^{-}(\hat c)/n\le k^{-}(\hat c)/n=q^{-}(\hat c)$
and $\tilde q^{+}(\hat c)=\tilde k^{+}(\hat c)/n\ge k^{+}(\hat c)/n
=q^{+}(\hat c)$, using $n\ge1$; for $\hat c\notin\Chatcon$ both pairs of
bounds are $0$ and $1$. So $\tilde{\mathbf{q}}^{\mp}$ satisfy the hypothesis,
and (i)--(iii) apply.

\emph{The canonical box form.} Every $\tilde q^{\mp}(\hat c)$ is an integer
divided by $n$, or $0$ or $1$, hence rational. For $\hat c\in\Chatcon$,
\eqref{eq:ordering} of Lemma~\ref{lem:guard}\,(i) gives
$k^{-}(\hat c)\le k^{+}(\hat c)$, so
$0\le\tilde k^{-}(\hat c)\le k^{-}(\hat c)\le k^{+}(\hat c)\le
\tilde k^{+}(\hat c)$ and $0\le\tilde q^{-}(\hat c)\le\tilde q^{+}(\hat c)$;
for $\hat c\notin\Chatcon$, $0\le0\le1$. These are the conditions of the
canonical box form, and its denotation \eqref{eq:boxset} with these bounds is
\eqref{eq:widened-set}.

\emph{The tails.} By Lemma~\ref{lem:guard}\,(ii) the map
$k\mapsto\Pr[\Ncount{\hat c}^{-}<k]$ is non-decreasing and the map
$k\mapsto\Pr[\Ncount{\hat c}^{+}>k]$ is non-increasing on the integers, so
$\tilde k^{-}(\hat c)\le k^{-}(\hat c)$ and
$\tilde k^{+}(\hat c)\ge k^{+}(\hat c)$ give the first inequality of each
half of \eqref{eq:widened-tails}; the second is \eqref{eq:budget-met} of the
same lemma. Summing the first inequalities over $\hat c\in\Chatcon$ compares
the two tail sums, and summing the second ones bounds the sum at $k^{\mp}$ by
$\sum_{\hat c\in\Chatcon}(\varepsilon^{-}(\hat c)+\varepsilon^{+}(\hat c))$,
which the allocation of an admissible input keeps at most $\ecomp$
(Definition~\ref{def:endpoint-rule}\,(vi)).

\emph{(v).} With the degenerate envelope at $\tfrac12$, the two counts
$\Ncount{\hat c_0}^{-}$ and $\Ncount{\hat c_0}^{+}$ have one law; write
$\Ncount{}$ for a count binomial with parameters $(3,\tfrac12)$, whose law is
$(\tfrac18,\tfrac38,\tfrac38,\tfrac18)$ on $\{0,1,2,3\}$. The exact rule
takes \eqref{eq:kappaminus}: $\Pr[\Ncount{}<0]=0$ and
$\Pr[\Ncount{}<1]=\tfrac18$ are at most $\tfrac25$, while
$\Pr[\Ncount{}<k]\ge\Pr[\Ncount{}<2]=\tfrac12>\tfrac25$ for every $k\ge2$, so
$\kappa^{-}(\hat c_0)=1$. It takes \eqref{eq:kappaplus}:
$\Pr[\Ncount{}>2]=\tfrac18\le\tfrac25$ and $\Pr[\Ncount{}>3]=0$, while
$\Pr[\Ncount{}>k]\ge\Pr[\Ncount{}>1]=\tfrac12>\tfrac25$ for every $k\le1$, so
$\kappa^{+}(\hat c_0)=2$. The guard \eqref{eq:guard} is inert, since
$\lfloor\tfrac32\rfloor=1$ and $\lceil\tfrac32\rceil=2$, so
$k^{-}(\hat c_0)=1$ and $k^{+}(\hat c_0)=2$, and the tails at them are
$\Pr[\Ncount{}<1]+\Pr[\Ncount{}>2]=\tfrac18+\tfrac18=\tfrac14$. The replaced
tails are $\Pr[\Ncount{}<2]=\tfrac12$ and $\Pr[\Ncount{}>1]=\tfrac12$, each
above $\tfrac25$.

For the failure probabilities, let $\Chat_1,\Chat_2,\Chat_3$ be independent
with law $\honest$. The count of $\hat c_0$ in $\Chat_1^3$ is then binomial
with parameters $(3,\tfrac12)$, so it has the law of $\Ncount{}$. On a
two-letter alphabet a vector of $\simplex{\Chatset}$ is fixed by its
coordinate at $\hat c_0$, and the bounds at $\hat c_1$ are $0$ and $1$ by
\eqref{eq:boxbounds}, so a frequency vector lies in a set \eqref{eq:accset}
built from endpoints $(a,b)$ at $\hat c_0$ exactly when the count of
$\hat c_0$ lies in $\{a,\dots,b\}$. For $\acSet$, with $(a,b)=(1,2)$, the
failure probability is $\Pr[\Ncount{}\in\{0,3\}]=\tfrac14$. With $k^{-}$
replaced by $2$, $(a,b)=(2,2)$ and the failure probability is
$1-\Pr[\Ncount{}=2]=\tfrac58$; with $k^{+}$ replaced by $1$, $(a,b)=(1,1)$
and it is $1-\Pr[\Ncount{}=1]=\tfrac58$. Both exceed $\tfrac14$.
\end{proof}
